\documentclass[11pt]{article}

% Change "review" to "final" to generate the camera-ready version.
\usepackage[review]{acl}

\usepackage{times}
\usepackage{latexsym}
\usepackage[T1]{fontenc}
\usepackage[utf8]{inputenc}
\usepackage{microtype}
\usepackage{inconsolata}
\usepackage{graphicx}
\usepackage{amsmath}

% [PLACEHOLDER: paper title — method name undecided]
\title{Learning When to Act: Speculative Tool Execution from the Hidden States of Reasoning Models}

\author{Anonymous ACL submission}

\begin{document}
\maketitle

\begin{abstract}
% [PLACEHOLDER: abstract — write last, after experiments section stabilizes. <=200 words.]
Reasoning models pay for every tool call twice: wall-clock time while the tool
runs, and thinking tokens spent deriving a decision whose outcome is often
predictable. We present [METHOD-NAME], a speculator that reads the frozen main
model's hidden states and predicts upcoming tool calls, their arguments, and
the benefit of acting early. [PLACEHOLDER: headline numbers.]
\end{abstract}

\section{Introduction}

Reasoning models use external tools by interleaving long chains of thought
with tool calls. The call itself is synchronous: the model thinks for hundreds
of tokens, commits to a call, and then waits for the result before thinking
further. This design bills every call twice. The user pays wall-clock time
while the tool runs, and pays again for the thinking tokens the model spends
deriving a decision whose outcome is often predictable, especially when an
agent serves many similar requests over its deployment lifetime.

Speculative execution is the natural answer, and recent systems have started
to apply it to tool calls. SPORK~\citep{spork2026} forks the model to act as
its own predictor of the upcoming call; PASTE~\citep{paste2026} mines
control-flow patterns from past workflows and pre-executes the calls they
imply; SIA~\citep{sia2026} speculatively executes calls whose arguments may
still be incomplete. These systems differ in what they predict, but they share
one design choice: the decision of \emph{when} to act is driven by output-side
confidence or by workflow structure, never by the state of the model that is
doing the thinking. Our measurements show this choice is not a detail.
Injecting a tool result 25--50 tokens before the model would issue the call
yields \emph{negative} returns: the model has entered the decision-critical
span of its reasoning, and the interruption costs more than the head start
saves.
% [PLACEHOLDER: dead-zone numbers] % source: run_id=hotpot_inject 第二波
We call this span the dead zone. Confidence-triggered firing lands in it
systematically, because prediction confidence peaks exactly where the model is
about to commit.
% [PLACEHOLDER: closed-loop threshold sweep numbers] % source: run_id=closed_loop_v0
The benefit of speculation is concentrated where confidence is lowest: very
early in the thinking, far ahead of the call point.

This paper argues that the missing signal lives in the model's hidden states,
and that it can be learned. We present [METHOD-NAME], a speculator: a set of
lightweight heads that read the frozen main model's hidden states at every
decoding step and predict, jointly, whether a call is coming, which tool it
will use, what its arguments will be, and how much is gained by acting now.
The argument head follows the shape of the argument space: calls with no
arguments are speculated most aggressively; calls whose arguments come from a
closed candidate set visible in the context are checked against that set
before speculation; free-form arguments are speculated conservatively or not
at all.
% [PLACEHOLDER: three-tier share table] % source: run_id=router_stats
The timing head is trained on labels produced by oracle replay: we re-run
stored trajectories, inject the known-correct result at varying lead times,
and record the realized benefit, so the head learns the dead-zone geometry
directly instead of inheriting it from confidence.

The speculator deploys in two modes with separate accounting. In prefetch
mode, a triggered call runs in the background while thinking continues
unchanged; the token stream is untouched, accuracy is preserved by
construction, and the tool's latency is hidden behind decoding. In truncate
mode, a triggered call ends the thinking immediately; this saves thinking
tokens and trades a controlled fraction of wrong calls, tunable by the firing
threshold $\theta$. On [BENCHMARK], truncate mode at $\theta=0.8$ saves
[PLACEHOLDER]\% of thinking tokens with [PLACEHOLDER] accuracy change,
% source: run_id=bfcl_v2fix_replay
and prefetch mode hides [PLACEHOLDER]\% of tool latency at zero accuracy
cost.

Our contributions: (1) a measurement of the injection-timing curve showing a
dead zone before the call point, which output-side confidence cannot avoid;
(2) [METHOD-NAME], a hidden-state speculator with a learned timing head and a
tiered argument head; (3) an evaluation protocol that reports speed and
accuracy jointly, as savings curves against accuracy-loss curves rather than
prediction accuracy alone; (4) [PLACEHOLDER: benchmark / learning-curve
contribution, pending merged-track design].
% 中文备注：贡献 (4) 待合体版 benchmark 定型

\section{Related Work}

\paragraph{Speculative execution of tool calls.}
SPORK~\citep{spork2026} needs no training: it forks the model mid-thought and
lets the fork race ahead to predict the full call, reaching
[PLACEHOLDER: 74.6--99.6\%] tool-name accuracy, at the price of a second
decoding stream. PASTE~\citep{paste2026} and B-PASTE~\citep{bpaste2026}
predict from mined workflow regularities and schedule speculation by
estimated critical-path reduction; the estimate comes from workflow
structure, so it cannot see effects that exist only in the model's reasoning,
the dead zone among them. Cost-aware schedulers~\citep{costaware2026} price
speculation in dollars with an expected-value rule, again without reading the
model. Speculate-with-Memory~\citep{swm2026} adds an external memory to its
predictor but still treats the main model as a black box. Against this line,
our speculator reads the main model's hidden states, adds no second decoding
stream, and learns timing from realized benefit rather than structure or
confidence.

\paragraph{Reading internal states for tool decisions.}
When2Tool~\citep{when2tool2026} shows that whether a tool is needed is
linearly decodable from hidden states (AUROC 0.89--0.96), but reads once,
before generation, and stops at the binary decision.
PPRAG~\citep{pprag2026} attaches a predictor to middle-layer states and
prefetches retrievals mid-decoding, cutting end-to-end latency by 43.5\%;
retrieval is its only tool, and a query its only argument. A recent probe
study~\citep{depgraph2026} decodes the dependency topology between future
calls from the residual stream and states explicitly that it makes no
behavioral-control claims. OLIVIA~\citep{olivia2026} runs a contextual bandit
over frozen hidden states to pick among candidate actions in repeated
deployments, without speculation or prefetching. These works establish that
the signals we need exist; none closes the loop from signal to speculative
execution of general tool calls with arguments, timing, and post-hoc
verification.

\paragraph{Asynchronous execution and mid-generation injection.}
AsyncLM~\citep{asynclm2024} removes the wait: calls run asynchronously and an
interrupt protocol notifies the model in-flight, cutting task latency by up
to 5.4$\times$. The call is still authored by the model itself; AsyncLM
removes waiting, not thinking. SIA~\citep{sia2026} injects finished results
into the context following the same mechanism. IdleSpec~\citep{idlespec2026}
fills tool-wait idle time with candidate plans, trading idle compute for
accuracy where we trade prediction for latency. Ghost-call privacy
analysis~\citep{ghost2026} shows that speculative calls leak at issue time,
which motivates the read-only restriction in our prefetch mode. Our
speculator composes with this line: a speculated call can be dispatched on an
asynchronous substrate, and our verification step reconciles it with the call
the model eventually issues.

\paragraph{Agents on repeated tasks.}
% [PLACEHOLDER: 此段待 benchmark 部分定型后重写，当前仅占位。]
AgentCL~\citep{agentcl2026} motivates wasted inference on recurring tasks
without measuring it; caching layers reuse results but not decisions;
Speculate-with-Memory~\citep{swm2026} draws a learning curve in estimated
latency rather than wall-clock time.

\section{Method}
% [PLACEHOLDER: 四头结构 / 三档参数产线 / oracle 回放标签 / 两种部署模式]

\section{Evaluation Protocol: Joint Speed--Accuracy Learning Curves}
\label{sec:protocol}
% 本节即 C3 协议节(T15 起笔,2026-07-30)。设计依据:
% benchmark_design/L3L4_and_metrics_draft.md;数字溯源见各处 source 注释。

Existing streaming benchmarks measure whether agents get more accurate with
experience; the cost of that experience is reported incidentally, if at all.
Our protocol makes amortized cost a first-class axis: every reported accuracy
number is paired with the cumulative tokens spent reaching it, on task streams
whose similarity structure is controlled. The protocol answers three
questions: how much does memory save as tasks repeat, what does competence
cost, and when does remembered experience actively hurt.

\subsection{Task Streams and the Similarity Ladder}
\label{sec:ladder}

A benchmark instance is a stream of $K$ episodes. Every subject runs each
stream twice under identical seeds and task instances: a \textsc{mem} arm
with its memory mechanism enabled, and a \textsc{nomem} arm with memory
disabled. All comparisons are paired at the (run, position, task instance)
level; the \textsc{nomem} arm absorbs the difficulty fluctuations of the
stream so that differences read out the memory mechanism alone.

Streams are drawn at five similarity levels, numbered by the degree to which
memory \emph{should} help. L0 streams contain unrelated tasks
(pairwise similarity $<0.3$): a well-behaved memory should be inert here, and
any accuracy drop on L0 measures the mechanism's indiscriminate overhead. L1
streams contain partially similar tasks (same task type, different objects or
entities). L2 streams contain repeats, from same-parameters-different-scene
up to exact repetition of one instance. L3 and L4 are structured probes
rather than points on the monotone axis. An L3 task is composed entirely of
components that each appeared in the stream's history while the combination
itself is new; it separates memories that store trajectories from memories
that store reusable skills. An L4 task is a trap: superficially as similar to
a seen task as a genuine repeat, but constructed so that the stored solution
fails. Admission to L4 requires both conditions to hold by construction --
embedding similarity to the seed task at or above the median similarity of
genuine L2 repeat pairs, and a mechanical replay of the seed's solution that
verifiably fails on the new task. Each admitted trap records its divergence
depth $d^\ast \in [0,1]$, the fraction of the replayed trajectory executed
before the first contradiction signal, and is tiered accordingly: T1
($d^\ast<0.25$, fails on contact), T2 (fails midway), and T3 (replay
completes without any contradiction signal and only the final grade fails).
T3 is the dangerous tier: the failure is silent in every process metric.

The ladder is not a synthetic worry. In a 357k-turn corpus of real coding
agent sessions, 68.4\% of adjacent-session pairs have tool-composition
similarity above 0.8, another 17.2\% fall between 0.3 and 0.8, and 96.1\% of
sessions have a near-duplicate predecessor within a 50-session window --
repeated similar work is the dominant regime of real agent workloads, not an
assumption. Among high-similarity pairs, 10.8\% cross project boundaries:
naturally occurring look-alikes with different ground truth, the L4
phenomenon in the wild.
% source: run_id=20260727_tracelab_simv0 (待 A 线 record 登记,数字见
% tracelab_analysis/RESULTS_v0.md;相似度仪器只见工具构成,边界已声明)

\subsection{Metrics}
\label{sec:metrics}

Let $u_{r,k}\in[0,1]$ and $c_{r,k}>0$ be the grade and cost of episode $k$ in
run $r$; cost is input plus output tokens, the only currency comparable
across serving hardware. Wall-clock is reported within one hardware
configuration only. $P_k$ denotes the paired grid at position $k$: runs whose
two arms both completed the full stream and both graded \emph{ok} at $k$.

\paragraph{Speedup@k.} The savings curve is the ratio of sums over the
paired grid,
\begin{equation}
\mathrm{Sp}@k \;=\; 1-\frac{\sum_{r\in P_k} c^{\textsc{mem}}_{r,k}}
                          {\sum_{r\in P_k} c^{\textsc{nomem}}_{r,k}},
\label{eq:spk}
\end{equation}
not the mean of per-run ratios, which is unstable under the heavy-tailed
per-episode savings we observe in practice. $\mathrm{Sp}@k$ is reported only
together with $\mathrm{Acc}@k=\tfrac{1}{|P_k|}\sum_{r\in P_k}
u^{\textsc{mem}}_{r,k}$ on the same grid: an agent that fails fast scores
spectacular savings, so a savings number without its paired accuracy is
meaningless. Positions with $|P_k|$ below a preregistered floor are left as
gaps, never filled by surviving runs.

\paragraph{Cost-to-competence.} Competence at position $k$ is the trailing
window mean $\hat p_{r,k}=\tfrac1w\sum_{j=k-w+1}^{k}u_{r,j}$ ($w{=}3$); the
attainment position is $k^\ast_r(\theta)=\min\{k\ge w:\hat
p_{r,k}\ge\theta\}$. The tuition to competence $\theta$ is
\begin{equation}
\mathrm{CTC}_r(\theta)\;=\;\textstyle\sum_{j=1}^{\min(k^\ast_r,K)} c_{r,j},
\label{eq:ctc}
\end{equation}
a sum that deliberately includes failed episodes: tuition includes the
retakes. Runs that never attain $\theta$ are censored at the full-stream
budget and reported with a $\ge$ sign, alongside the attainment rate
$\rho(\theta)$; averaging over attaining runs only is the textbook form of
survivorship bias and is excluded by construction.

\paragraph{Negative-transfer index.} On the L4 probe segment (seeding
episodes excluded), with both arms paired and ok-graded,
\begin{align}
\mathrm{NTI}_{\mathrm{acc}}&=\bar u^{\textsc{nomem}}-\bar u^{\textsc{mem}},
\label{eq:nti-acc}\\
\mathrm{NTI}_{\mathrm{cost}}&=
\frac{\bar c^{\textsc{mem}}-\bar c^{\textsc{nomem}}}{\bar c^{\textsc{nomem}}},
\label{eq:nti-cost}
\end{align}
both signed so that positive means memory hurt. The two components are never
merged: faster-but-wrong (T3 signature) and slower-and-wrong (T2 signature)
are different failure modes that a weighted sum would average away. Results
are stratified by trap tier, yielding a dose--response curve, and accompanied
by the misuse rate: the fraction of probe episodes whose action prefix copies
the stored solution past its divergence depth. A significant
$\mathrm{NTI}_{\mathrm{acc}}$ with near-zero misuse rate indicates the drop
has another cause and may not be attributed to experience misuse.

\subsection{Aggregation Rules}
\label{sec:rules}

Five rules, implemented as assertions in the released analysis code rather
than as conventions. (1) Per-position statistics use only runs that completed
all $K$ episodes. (2) Every \textsc{mem}--\textsc{nomem} comparison is paired
on run, position, and task instance; a cell missing either arm is dropped
whole. (3) Headline metrics are computed on ok-graded episodes, and the
ok-rate of every condition is itself reported: a condition that inflates its
ok-only mean by shedding hard episodes is exposed by its ok-rate. (4)
Censoring is explicit: no reachers-only means, ever. (5) Wall-clock never
aggregates across hardware; tokens are the cross-machine currency.

\subsection{Baselines}
\label{sec:baselines}

Every subject is bracketed by three memoryless references. The
\textsc{nomem} floor is the paired arm defined above. The oracle-replay
ceiling replays the cheapest known successful solution of an instance at
zero token cost once one exists in the stream's history; no memory system
can beat free replay. The ceiling is informative precisely because it is
narrow: verbatim replay is defined only where instances truly repeat, and on
our L2 streams it saves at most 21\% of stream tokens while leaving success
unchanged -- the expensive episodes are the failures, which replay cannot
fix. Subjects must therefore be read against the ceiling of their level: savings
beyond the ceiling can only come from generalization, not recall.
% source: run_id=20260730_oracle_ceiling_8bfull (待 A 线 record 登记,数字见
% benchmark_design/ORACLE_CEILING_8bfull.md,记录基座 run_id=20260728_fig1_8bfull)
The third reference stuffs the full interaction history into the context
window, a baseline prior work found surprisingly strong; it prices the
alternative of not having a memory mechanism at all, only a long context.

\section{Experiments}
% [PLACEHOLDER: benchmark 定型后填；主图=省 token 曲线 vs 准确率损失曲线]

\section{Conclusion}
% [PLACEHOLDER]

\section*{Limitations}
% [PLACEHOLDER: ACL 必需节。候选内容：合成任务边界 / 协议遵守对模型敏感 /
% 只测 Qwen 系 / 授权措辞的必要性按模型规模限定]
TODO.

\bibliography{references}

\appendix

\section{Appendix}
% [PLACEHOLDER]

\end{document}
